% Austrian Journal of Statistics article structure starter.
\documentclass[11pt,a4paper]{article}
\usepackage[margin=2.25cm]{geometry}
\usepackage{amsmath,booktabs,tabularx,xcolor,hyperref,microtype}
\hypersetup{colorlinks=true,urlcolor=blue!50!black,linkcolor=blue!50!black,citecolor=blue!50!black}
\setlength{\parindent}{0pt}
\setlength{\parskip}{0.55em}
\title{Estimating service reliability from incomplete observations}
\author{Clara Novak\\\small Institute for Data Analysis\\\small Example University}
\date{Austrian Journal of Statistics --- Volume 55, Issue 1 (2026)}
\begin{document}
\maketitle
\begin{abstract}
This article demonstrates an AJS-style manuscript with a concise summary,
author information, reproducible methods, and clearly labelled tables. The
fictional study estimates weekly service reliability when a small fraction of
observations is missing.
\end{abstract}
\textbf{Keywords:} reliability; missing data; bootstrap; reproducibility

\section{Introduction}
Operational datasets often contain gaps that are small in number but unevenly
distributed over time. A transparent estimator should report both the point
estimate and the sensitivity to plausible missingness mechanisms.

\section{Model}
Let $y_t$ be the observed weekly success rate and $m_t$ the fraction of
missing records. We estimate the adjusted rate as
\begin{equation}
  \widehat{r}_t = \frac{y_t}{1-m_t}, \qquad 0 \leq m_t < 1.
\end{equation}
The reported interval is obtained from 2,000 bootstrap resamples stratified
by week.

\section{Results}
Table~\ref{tab:estimate} gives the illustrative estimates. All values are
fictional and should be replaced with a reproducible analysis.
\begin{table}[ht]
\centering
\caption{Illustrative weekly reliability estimates.}
\label{tab:estimate}
\begin{tabularx}{0.78\linewidth}{@{}lrrX@{}}
\toprule
Week & Observed & Adjusted & 95\% interval \\
\midrule
1 & 0.91 & 0.94 & [0.90, 0.97] \\
2 & 0.88 & 0.92 & [0.87, 0.96] \\
3 & 0.93 & 0.95 & [0.92, 0.98] \\
\bottomrule
\end{tabularx}
\end{table}

\section{Discussion}
The adjustment is useful only when the missingness assumption is defensible.
Sensitivity analysis should therefore be included whenever the missing
records could be related to the outcome.

\section{Conclusion}
An explicit estimator, a stated resampling procedure, and a compact table
make the statistical argument easier to inspect and reproduce.

\section*{Address}
Clara Novak\\
Institute for Data Analysis, Example University\\
Example Street 12, 1010 Example City, Austria\\
\texttt{clara.novak@example.org}

\begin{thebibliography}{9}
\bibitem{efron}
Efron, B. and Tibshirani, R. J. (1993). \emph{An Introduction to the
Bootstrap}. Chapman and Hall.
\end{thebibliography}
\end{document}
